\documentclass[11pt,a4paper]{article}
\usepackage[T1]{fontenc}
\usepackage{isabelle,isabellesym}
\usepackage{amsmath,amssymb}
\usepackage{pdfsetup}

\urlstyle{rm}
\isabellestyle{it}

\begin{document}
\emergencystretch=3em

\title{ERC-TRUST: Mechanized Regulatory Execution Semantics}
\author{Jinwook Kim (Jay)}
\maketitle

\begin{abstract}
This Isabelle/HOL session mechanizes the abstract regulatory-action semantics
used by ERC-TRUST. It distinguishes six RCP legal-effect actions from the
seven transition labels of the inherited regulatory state machine, specifies
typed authorization and observable outcomes, and checks an executable
manifest kernel against the abstract model. The session depends on
\texttt{Cross\_Domain\_State\_Preservation} as a parent Isabelle session.
The result is model-level evidence only. It does not establish legal title,
judicial validity, off-chain evidence truth, Solidity-bytecode conformance,
or deployment correctness.
\end{abstract}

\tableofcontents

\section{Session Boundary}

The parent session supplies the cross-domain regulatory state machine and the
regulatory-action composition theory. ERC-TRUST imports those theories as
formal dependencies without copying their source into this entry. The
ERC-TRUST session owns only its regulatory execution semantics, compatibility
boundary, privileged-governance model, executable kernel, simulation
witnesses, claim boundary, and proof-trust audit.

\section{Normative Context}

ERC-TRUST is designed as a thin conformance extension over ERC-20 and
ERC-7943, with regulatory action meaning supplied by ERC-8319
\cite{ERC7943,ERC8319}. These specifications provide the normative context;
the Isabelle theories define and verify the declared abstract execution
domain.

\section{Verification Boundary}

The theories prove internal consistency and stated safety properties over the
declared model. They deliberately do not claim that an \texttt{Applied}
outcome proves lawful forfeiture, actual sale completion, debt discharge,
rightful ownership, or truth of an external policy or identity assertion.
Implementation-level claims require separately bound Solidity, Certora, and
Kontrol evidence.

\input{session}

\bibliographystyle{plain}
\bibliography{root}

\end{document}